\chapter{Paires et paniers}\label{ch:s1:pairs-and-baskets}

Deux actions synthétiques avaient évolué ensemble pendant un an, leurs prix normalisés n'étaient séparés que d'un écart-type de 2.4~\%~: elles avaient été conçues comme
des substituts proches, et la \omterm{def:s1:pairs-and-baskets:distance}{méthode des distances} les a retenues. La paire s'ouvrit lorsque leur spread atteignit deux écarts-types, se referma, puis se rouvrit. Puis, le 111e jour de trading, l'une des sociétés fut rachetée avec une prime. Le spread bondit à 35~\%, la jambe vendeuse cessa d'être négociée, et la paire
ne se referma jamais~: elle a perdu 24.6~\% par dollar et par jambe. Le \omterm{def:s1:pairs-and-baskets:pairs}{pairs trading} est l'histoire la plus ancienne de l'\omterm{def:s1:anatomy-of-a-stat-arb-book:statarb}{arbitrage statistique} et la plus simple à expliquer~: trouver deux actions
qui évoluent ensemble, et miser sur leur retour l'une vers l'autre lorsqu'elles s'écartent. Ce chapitre traite de la manière dont les deux actions sont trouvées, et pourquoi la plupart des méthodes
trouvent les mauvaises, et ce qui arrive lorsque le récit cesse d'être vrai. La construction est \texttt{firm.pairsel}.

\section{Sélection~: distance, cointégration, copules}

\begin{definition}[Pairs trading]\label{def:s1:pairs-and-baskets:pairs}
Le \emph{pairs trading}\index{pairs trading} est la stratégie qui consiste à sélectionner deux titres dont les cours ont évolué ensemble, et à prendre une position longue sur celui qui a pris du retard
et une position courte sur l'autre lorsque leurs cours divergent, positions à dénouer lorsqu'ils convergent.
\end{definition}

\begin{definition}[Méthode des distances, période de formation, période de trading]\label{def:s1:pairs-and-baskets:distance}
La \emph{méthode des distances}\index{méthode des distances} sélectionne les paires d'actions dont les trajectoires de prix normalisés (indices de rendement total cumulé partant de
un) dont la somme des carrés des écarts est la plus petite sur une \emph{période de formation}\index{période de formation}, et les négocie pendant la
\emph{période de trading}\index{période de trading}, en ouvrant la paire lorsque son spread dépasse un multiple de son écart-type sur la période de formation, et en la clôturant
lorsque les cours se croisent.
\end{definition}

Gatev, Goetzmann et Rouwenhorst ont donné à la \omterm{def:s1:pairs-and-baskets:distance}{méthode des distances} sa forme standard~: douze mois de formation, six mois de trading, ouverture à deux
écarts-types historiques, clôture au prochain croisement ou à la fin des six mois~; les 20 premières paires. Sur des données quotidiennes américaines de 1962 à 2002, elle a rapporté des rendements annualisés
excédentaires pouvant atteindre 11~\% sur des portefeuilles autofinancés de paires, généralement supérieurs aux estimations prudentes des coûts. La méthode de cointégration (livre~4,
chapitre~20) remplace la distance par un test d'Engle--Granger~: on régresse un log-prix sur l'autre et l'on teste les résidus pour une racine unitaire. La méthode des copules
modélise les rendements quotidiens de la paire à l'aide d'une copule (livre~4, chapitre~15) et prend position lorsqu'un rendement est invraisemblablement élevé ou faible compte tenu de l'autre. L'enquête de Krauss
recense plus d'une centaine d'articles répartis en cinq familles, et Rad, Low et Faff ont comparé les trois principales sur les actions américaines de 1962 à 2014~: rendements mensuels excédentaires moyens
de 91, 85 et 43 points de base avant coûts, pour la \omterm{def:s1:pairs-and-baskets:distance}{méthode des distances}, la cointégration et la méthode des copules, et de 38, 33 et 5 après coûts.

Le marché synthétique ne contient pour l'instant aucune paire à trouver~: ses actions partagent des facteurs, mais aucune n'est le substitut d'une autre, si bien que chaque spread s'éloigne au gré de la
différence de deux marches aléatoires. \texttt{MarketConfig.twin\_share} (nouveau dans ce chapitre, désactivé par défaut) apparie 10~\% des titres initialement cotés en substituts
proches, soit cinquante paires dont le spread des log-prix est un AR(1) stationnaire, de temps de retour à la moyenne de dix jours et d'écart-type de 2~\%. Tout le reste
concernant les jumeaux est ordinaire~: ils peuvent être rachetés, radiés ou fractionnés comme n'importe quelle action. Comme les paires sont connues, le chapitre peut compter ce que
chaque méthode trouve.

\section{Règles de trading et fenêtre de formation}

Sur le marché des jumeaux, seize cycles de douze mois de formation et six mois de trading, de l'année 3 à l'année 10, avec la règle de Gatev--Goetzmann--Rouwenhorst,
1~\$ par jambe et cinq points de base par unité échangée~:

\begin{center}
\small
\setlength{\tabcolsep}{4pt}%
\begin{tabular}{@{}p{5.0cm}ccccc@{}}
\toprule
sélection & distance & copule & deux étapes & Engle--Granger & \omterm{def:s1:pairs-and-baskets:twin}{jumeau synthétique}\\
\midrule
paires implantées parmi les 320 sélectionnées & 311 & 311 & 113 & 40 & ---\\
ratio de Sharpe avant coûts & 5.34 & 4.39 & 1.16 & $-0.09$ & $-0.45$\\
ratio de Sharpe après coûts & 5.05 & 3.89 & 1.03 & $-0.17$ & $-0.49$\\
rendement net annuel (1~\$ par jambe) & 12.2~\% & 11.6~\% & 5.9~\% & $-1.4~\%$ & $-2.6~\%$\\
ouvertures par paire et par cycle & 2.11 & 4.28 & 2.24 & 2.02 & 0.87\\
paires encore ouvertes à la fin & 51~\% & 89~\% & 75~\% & 81~\% & 61~\%\\
\midrule
ratio de Sharpe après coûts, sans paires implantées & $-0.43$ & $-0.50$ & $-0.11$ & $-0.47$ & $-0.35$\\
\bottomrule
\end{tabular}
\end{center}

La \omterm{def:s1:pairs-and-baskets:distance}{méthode des distances} trouve les substituts~: 311 de ses 320 paires sont des jumeaux implantés, et elle dégage un ratio de Sharpe de 5.05. Sur le marché par défaut, où il
n'existe aucune, la même méthode perd~: les paires qu'elle retient sont celles dont les trajectoires sont restées proches pendant un an, généralement deux actions à faible volatilité, et
84~\% d'entre elles sont encore ouvertes à la fin de leurs six mois, car rien ne les ramène. Gatev, Goetzmann et Rouwenhorst ont relevé le même
effet de sélection~: en moyenne 82~\% des actions de leurs vingt premières paires étaient des sociétés de services publics. L'indice des copules sur les mêmes paires négocie deux fois plus souvent et rapporte un peu
moins (son indice de mauvaise évaluation dérive, si bien que les paires restent ouvertes). La pire colonne du tableau est celle qui paraît la plus rigoureuse.

\begin{omfigure}
\begin{tikzpicture}
\begin{axis}[width=11cm,height=5.0cm,xlabel={\small jour de la période de trading},ylabel={\small spread~; P\&L par jambe de 1~\$},tick label style={font=\footnotesize},
  xmin=0,xmax=126,ymin=-0.3,ymax=0.4,grid=major,grid style={gray!25},table/col sep=comma,legend pos=north west,legend style={font=\scriptsize}]
\addplot[omDef,very thick] table[x=day,y=spread]{figdata/strategies-1/04-pairs-and-baskets/takeover.csv};
\addplot[omThm,thick] table[x=day,y=pnl]{figdata/strategies-1/04-pairs-and-baskets/takeover.csv};
\addplot[black!50,dashed] coordinates {(0,0.048) (126,0.048)};
\addplot[black!50,dashed] coordinates {(0,-0.048) (126,-0.048)};
\addplot[black!60,dotted,thick] coordinates {(111,-0.3) (111,0.4)};
\legend{spread,P\&L cumulé,$\pm 2$ écarts-types de formation}
\end{axis}
\end{tikzpicture}

\omcaption{L'exemple d'ouverture du chapitre~: deux substituts implantés sur le marché synthétique pendant leur \omterm{def:s1:pairs-and-baskets:distance}{période de trading} de six mois. Le spread (différence des
prix normalisés) ouvre la paire deux fois~; le jour 111 (en pointillés), le titre vendu par la paire est racheté avec une prime, et la paire ne se referme jamais. Données~:
\texttt{s1\_pairs.takeover}.}\label{fig:s1:pairs-and-baskets:takeover}
\end{omfigure}

\section{Pourquoi les tests de cointégration échouent à grande échelle}

Classer les paires selon leurs p-values d'Engle--Granger devrait faire ressortir les paires implantées, puisque ce sont les seules cointégrées. Il en trouve 40 sur 320. La raison en est
les tests multiples (livre~4, chapitre~12) avec un test dont les p-values sont fausses. Chaque cycle teste chaque paire d'actions d'un même secteur, soit environ 46\,600 paires~;
sur seize cycles, soit 745\,465 tests. Les p-values proviennent de la loi de la statistique d'Engle--Granger pour deux marches aléatoires gaussiennes indépendantes, simulée
une seule fois (\texttt{eg\_null}). Les rendements du marché synthétique ont des queues épaisses et une volatilité qui se regroupe, comme dans les données réelles, et dans ces conditions le test
rejette trop souvent~: sur le marché par défaut, où aucune paire n'est cointégrée, 7.4~\% des paires ont des p-values inférieures à 5~\%, et 0.32~\% inférieures à 0.1~\%, soit trois fois le
taux nominal (la \cref{fig:s1:pairs-and-baskets:calib}). Les plus petites p-values sont dominées par ces faux rejets, et non par les paires implantées.

\begin{omfigure}
\begin{tikzpicture}
\begin{axis}[width=7.2cm,height=5.6cm,xmode=log,ymode=log,xlabel={\small niveau nominal},ylabel={\small part des paires rejetées},
  tick label style={font=\footnotesize},xmin=0.0007,xmax=0.15,ymin=0.0007,ymax=0.2,grid=major,grid style={gray!25},legend pos=north west,
  legend style={font=\scriptsize},table/col sep=comma]
\addplot[omThm,very thick,mark=square*] table[x=nominal,y=actual]{figdata/strategies-1/04-pairs-and-baskets/calibration.csv};
\addplot[black!50,dashed] coordinates {(0.0007,0.0007) (0.15,0.15)};
\legend{p-values d'Engle--Granger,un test valide}
\end{axis}
\end{tikzpicture}

\omcaption{Calibrage du test d'Engle--Granger avec une hypothèse nulle de marche aléatoire gaussienne sur le marché synthétique sans paires implantées (742\,028 paires d'un même secteur
testées)~: la part des paires dont la p-value est inférieure à chaque niveau, en regard de ce niveau. Données~: \texttt{s1\_pairs.calibration}.}\label{fig:s1:pairs-and-baskets:calib}
\end{omfigure}

Sur le marché des jumeaux, les comptages sont encore plus nets. Sur 58\,024 paires rejetées à 5~\% (7.8~\% de celles testées), 530 sont implantées, soit 84.7~\% des 626 paires implantées
que les cycles pouvaient tester~: le test a de la puissance. La procédure de Benjamini et Hochberg, réglée sur un taux de fausses découvertes (FDR) de 5~\%, retient 1\,139 paires, dont 32 sont
implantées~: 97.2~\% de ses découvertes sont fausses. La procédure est correcte~; son entrée ne l'est pas, et un contrôle des fausses découvertes fondé sur des p-values mal calibrées
ne contrôle rien. La persistance dit la même chose, vue de l'autre côté~: parmi les paires rejetées à 5~\% pendant une année de formation et toujours cotées un an plus tard,
7.9~\% sont rejetées à nouveau l'année suivante, à peine au-dessus du taux de faux rejets du test. Les paires implantées sont rejetées à nouveau dans 84.6~\% des cas,
les autres dans 7.2~\%.

La méthode en deux étapes, une réparation courante, présélectionne les 200 paires de plus petite distance puis les classe selon la p-value d'Engle--Granger. Elle retient 113 paires implantées
et rapporte 1.03~: meilleure que le test seul, moins bonne que la distance seule, car à l'intérieur de la présélection le test continue de préférer ses faux rejets.

\section{Pourquoi les paires ont disparu, et ce qui les a remplacées}

Les données publiées indiquent que les profits ont diminué. Do et Faff ont confirmé une baisse continue de la rentabilité du \omterm{def:s1:pairs-and-baskets:pairs}{pairs trading}, la stratégie
restant solide pendant les périodes de turbulence prolongée, dont la crise de 2008~; Rad, Low et Faff ont constaté qu'à partir de 2009 les méthodes de la distance et de la cointégration ont trouvé
beaucoup moins d'opportunités de trading, tandis que celle des copules est restée stable. Le marché synthétique donne le mécanisme, sinon l'histoire~: la méthode de la distance
rapporte exactement autant que le marché contient de substituts proches, et rien d'autre. Une méthode qui sélectionne sur la co-variation passée des cours trouve
substituts là où ils existent, et des coïncidences de faible volatilité là où ils n'existent pas.

Ce qui a remplacé la paire classique, c'est le reste de cette partie du livre~: les résidus face à de nombreux facteurs (chapitre~3), où le «~partenaire~» est un modèle factoriel~;
la valeur relative structurelle (chapitre~12), où les substituts sont deux catégories d'actions, une société holding et ses participations, ou un fonds et ses actifs, et où le
lien est contractuel~; et les portefeuilles de milliers de petits paris (chapitre~1), où aucune paire isolée ne compte.

\section{Paniers~: une action contre son jumeau synthétique}

\begin{definition}[Jumeau synthétique]\label{def:s1:pairs-and-baskets:twin}
Le \emph{jumeau synthétique}\index{jumeau synthétique} d'une action est un portefeuille d'autres actions, généralement de son secteur, dont les rendements suivent ceux de l'action pendant une
\omterm{def:s1:pairs-and-baskets:distance}{période de formation}~; il est estimé par une régression des rendements de l'action sur ceux de ce portefeuille, et l'écart entre l'action et son jumeau se négocie comme celui d'une paire.
\end{definition}

Le panier généralise la paire~: au lieu d'un seul partenaire, trente pairs, pondérés par une régression ridge sur les rendements de l'année de formation. Sur le marché synthétique
il perd dans les deux versions ($-0.49$ et $-0.35$). Une régression sur trente pairs ajuste bien les mouvements passés de l'action liés aux facteurs, mais ses mouvements spécifiques lui appartiennent~:
ils sont indépendants de ceux de chacun des pairs, si bien que l'écart avec un panier quelconque est la marche aléatoire spécifique de l'action, et une marche aléatoire ne revient pas à la moyenne. Le panier
fonctionne là où l'action et ses pairs partagent quelque chose au-delà des facteurs (un client commun, un prix réglementé), ce qui est la question à poser sur tout
\omterm{def:s1:pairs-and-baskets:twin}{jumeau synthétique} avant de le négocier.

\begin{omfigure}
\begin{tikzpicture}
\begin{axis}[width=9.5cm,height=5.4cm,xlabel={\small année},ylabel={\small rendement net cumulé (\%, 1~\$ par jambe)},tick label style={font=\footnotesize},
  xmin=2,xmax=10,grid=major,grid style={gray!25},legend pos=outer north east,legend style={font=\scriptsize},table/col sep=comma]
\addplot[omDef,very thick] table[x=year,y=distance]{figdata/strategies-1/04-pairs-and-baskets/books.csv};
\addplot[omDef!60,thick] table[x=year,y=copula]{figdata/strategies-1/04-pairs-and-baskets/books.csv};
\addplot[omThm!60,thick] table[x=year,y=twostep]{figdata/strategies-1/04-pairs-and-baskets/books.csv};
\addplot[omThm,thick] table[x=year,y=eg]{figdata/strategies-1/04-pairs-and-baskets/books.csv};
\addplot[black!50,thick,dashed] table[x=year,y=twin]{figdata/strategies-1/04-pairs-and-baskets/books.csv};
\legend{distance,copule,deux étapes,Engle--Granger,jumeau synthétique}
\end{axis}
\end{tikzpicture}

\omcaption{Rendements nets cumulés des cinq méthodes de sélection sur le marché synthétique avec cinquante paires implantées de substituts. Données~:
\texttt{s1\_pairs.run}.}\label{fig:s1:pairs-and-baskets:books}
\end{omfigure}

\section{Fiches de stratégie}

\begin{strategyfile}{Paires de la méthode des distances}\label{strat:s1:pairs-and-baskets:distance}
\sfield{Qui paie, et pourquoi} Des flux qui écartent l'un de l'autre deux substituts proches~: le livre prête de la liquidité sur la paire.
\sfield{Instruments et plateformes} Actions liquides, achat de l'une et vente à découvert de l'autre~; mêmes horaires de bourse.
\sfield{Signal} L'écart des cours normalisés par son écart-type de formation~; ouverture à deux, clôture au croisement.
\sfield{Dimensionnement et exécution} 1~\$ par jambe et par paire, vingt paires~; les deux jambes négociées ensemble à la clôture.
\sfield{Coûts} Quatre unités négociées par aller-retour~; emprunt sur la jambe vendeuse.
\sfield{Comment elle meurt} Les paires sont des coïncidences, non des substituts~; une OPA ou une autre nouvelle brise une vraie paire (l'exemple d'ouverture du chapitre)~; trop de traders
sur les mêmes paires.
\sfield{Horizon, capacité, infrastructure} Quelques semaines à quelques mois~; faible capacité par paire~; un fichier de cours quotidiens et un traitement de formation.
\sfield{Backtest honnête} L'univers de formation tel qu'il était~; rendements de radiation et d'OPA dans les jambes~; paires conservées jusqu'à la fin de la \omterm{def:s1:pairs-and-baskets:distance}{période de trading}.
\sfield{Sources} Gatev, Goetzmann et Rouwenhorst (2006)~: jusqu'à 11~\% par an de rendements en excès, 1962--2002~; Do et Faff (2010)~: une baisse continue~; Rad,
Low et Faff (2016)~: 38 points de base par mois après coûts, 1962--2014.
\end{strategyfile}

\begin{strategyfile}{Paires de cointégration}\label{strat:s1:pairs-and-baskets:coint}
\sfield{Qui paie, et pourquoi} Comme pour les paires de la \omterm{def:s1:pairs-and-baskets:distance}{méthode des distances}, sélectionnées par un test d'écart stationnaire.
\sfield{Instruments et plateformes} Actions liquides, généralement du même secteur.
\sfield{Signal} Le résidu d'Engle--Granger d'un log-cours sur l'autre, en écarts-types de formation.
\sfield{Dimensionnement et exécution} Jambes dans le rapport du coefficient de cointégration~; ouverture à deux, clôture à zéro.
\sfield{Coûts} Comme pour les paires de la \omterm{def:s1:pairs-and-baskets:distance}{méthode des distances}.
\sfield{Comment elle meurt} Des p-values mal calibrées sur des dizaines de milliers de tests remplissent le livre de fausses paires~; des relations qui se rompent.
\sfield{Horizon, capacité, infrastructure} Quelques semaines à quelques mois~; un test par paire candidate et par formation.
\sfield{Backtest honnête} Compter les tests~; vérifier la calibration du test sur des données sans cointégration~; présélectionner selon l'économie, et non sur la seule p-value.
\sfield{Sources} Rad, Low et Faff (2016)~: 33 points de base par mois après coûts~; moins d'opportunités après 2009.
\end{strategyfile}

\begin{strategyfile}{Paires à copule}\label{strat:s1:pairs-and-baskets:copula}
\sfield{Qui paie, et pourquoi} Comme pour les paires de la \omterm{def:s1:pairs-and-baskets:distance}{méthode des distances}~; le signal mesure à quel point le mouvement conjoint d'aujourd'hui est improbable.
\sfield{Instruments et plateformes} Comme pour les paires de la \omterm{def:s1:pairs-and-baskets:distance}{méthode des distances}.
\sfield{Signal} Un indice de mauvaise tarification~: la somme cumulée de $P(U_a \le u_a \mid U_b = u_b) - \tfrac12$ sous une copule ajustée pendant la \omterm{def:s1:pairs-and-baskets:distance}{période de formation}.
\sfield{Dimensionnement et exécution} Ouverture à un indice de $\pm 0.5$, clôture lorsqu'il change de signe.
\sfield{Coûts} Négocie plus souvent que la règle de la distance.
\sfield{Comment elle meurt} Comme ci-dessus~; l'indice dérive, et les paires restent ouvertes.
\sfield{Horizon, capacité, infrastructure} Quelques semaines~; ajustements des marginales et de la copule par paire.
\sfield{Backtest honnête} Marginales et copule ajustées sur les seules données de formation.
\sfield{Sources} Rad, Low et Faff (2016)~: 43 points de base par mois avant coûts et 5 après~; opportunités de trading stables après 2009.
\end{strategyfile}

\begin{strategyfile}{Action contre son panier synthétique}\label{strat:s1:pairs-and-baskets:basket}
\sfield{Qui paie, et pourquoi} Comme pour les paires, avec un panier de pairs comme partenaire.
\sfield{Instruments et plateformes} Une action et un panier de pairs de son secteur.
\sfield{Signal} Le résidu cumulé de l'action face à son jumeau de régression ridge, en écarts-types de formation.
\sfield{Dimensionnement et exécution} 1~\$ sur l'action contre 1~\$ dans le panier.
\sfield{Coûts} Les transactions du panier s'ajoutent à celles de l'action.
\sfield{Comment elle meurt} Les mouvements spécifiques de l'action n'ont de contrepartie dans aucun panier~; l'écart est alors une marche aléatoire.
\sfield{Horizon, capacité, infrastructure} Quelques semaines~; une régression par action et par formation.
\sfield{Backtest honnête} Pondérations ajustées sur les seules données de formation~; comparer avec le \omterm{def:s1:short-term-reversal:residual}{retournement résiduel} propre à l'action (chapitre~3).
\sfield{Sources} Simulation de ce chapitre uniquement~; aucun historique de performance public vérifié.
\end{strategyfile}

\section{Tutoriel~: la paire qui ne s'est jamais refermée}\label{tut:s1:pairs-and-baskets:tut}

\begin{tutorial}
\textbf{Objectif.} Sélectionner des paires par la distance, par des tests d'Engle--Granger et par les deux~; les négocier~; compter les fausses découvertes. \textbf{État final~:} le tableau et
la \cref{fig:s1:pairs-and-baskets:calib}.
\begin{tutsteps}
\item \textbf{Les tests}~: Engle--Granger vectorisé sur de nombreuses paires, p-values issues d'une hypothèse nulle simulée, Benjamini--Hochberg.
\omcode{code/firm/pairsel/firm_pairsel.py}{55}{86}{La statistique d'Engle--Granger sur de nombreuses paires, son hypothèse nulle et le contrôle des fausses découvertes.}\label{lst:s1:pairs-and-baskets:eg}
\item \textbf{La règle}~: ouverture à deux écarts-types de formation, clôture au croisement, positions dérivant avec les rendements.
\omcode{code/firm/pairsel/firm_pairsel.py}{89}{116}{La règle de la distance sur une paire.}\label{lst:s1:pairs-and-baskets:rule}
\item \textbf{Exécuter} \texttt{run(share, method)} pour les cinq méthodes sur les deux marchés, \texttt{discoveries}, \texttt{calibration}, \texttt{takeover} et
\texttt{fig\_pairs.py}.
\end{tutsteps}
\textbf{Pour aller plus loin.} Remplacer l'hypothèse nulle gaussienne par un bootstrap des résidus propres à chaque paire et recompter~; élargir l'écart implanté à 5~\% et observer la
méthode de la distance perdre ses jumeaux au profit de coïncidences de faible volatilité~; fermer les paires dont une jambe fait l'objet d'une offre publique d'achat.
\end{tutorial}

\section{Construction~: sélection de paires}\label{bld:s1:pairs-and-baskets:pairsel}

\begin{build}
\bfield{Rôle} Sélection des paires par la distance, la cointégration et la copule, et leur règle de trading, avec la comptabilité des tests multiples qu'exige la sélection.
\bfield{Interface} \texttt{normalise(P)}, \texttt{top\_pairs(N, k)}, \texttt{eg\_tau(y, x)}, \texttt{eg\_null(T, reps, seed)}, \texttt{pvalue(tau, null)},
\texttt{bh(p, q)}, \texttt{trade\_pair(ra, rb, spread, sd, k, cost)}, \texttt{copula\_fit(ra, rb)}, \texttt{copula\_h(fit, ra, rb)}, \texttt{twin(y, X, lam)}.
\bfield{Règles} Données de formation seules utilisées pour la sélection et la mise à l'échelle~; chaque test est compté~; rendements de radiation et offres publiques d'achat intégrés dans les jambes.
\bfield{Tests d'acceptation} \texttt{code/firm/pairsel/tests/}~: la distance retrouve les quasi-copies implantées~; le quantile à 5~\% de l'hypothèse nulle proche de la valeur de MacKinnon, la puissance du test
sur un écart stationnaire et sa taille sur des marches aléatoires~; Benjamini--Hochberg sur un exemple de manuel~; le P\&L de la règle calculé à la main~; la corrélation de la copule et
les probabilités conditionnelles~; les poids ridge~; la loi normale inverse.
\bfield{Pour aller plus loin} Un bootstrap de l'hypothèse nulle par paire~; des tests de Johansen sur des paniers (\texttt{firm.coint} du livre~4)~; stop-loss et filtres d'offres publiques d'achat.
\end{build}

\begin{omsources}
\item E.~Gatev, W.~N.~Goetzmann and K.~G.~Rouwenhorst, ``Pairs trading: performance of a relative-value arbitrage rule'', \emph{Review of Financial Studies}
19(3), 2006 (NBER, document de travail 7032, 1999).
\item B.~Do and R.~Faff, ``Does simple pairs trading still work?'', \emph{Financial Analysts Journal} 66(4), 2010.
\item C.~Krauss, ``Statistical arbitrage pairs trading strategies: review and outlook'', \emph{Journal of Economic Surveys} 31(2), 2017.
\item H.~Rad, R.~K.~Y.~Low and R.~Faff, ``The profitability of pairs trading strategies: distance, cointegration and copula methods'', \emph{Quantitative
Finance} 16(10), 2016.
\end{omsources}

\section{Exercices}

\begin{exercise}[$\star$]\label{exo:s1:pairs-and-baskets:1}
Dix secteurs de 45 actions chacun~: combien y a-t-il de paires intra-sectorielles, et combien un test valide à 5~\% en rejetterait-il si aucune n'est cointégrée~?
\end{exercise}

\begin{exercise}[$\star$]\label{exo:s1:pairs-and-baskets:2}
Benjamini--Hochberg avec $q = 0.05$ et 1\,000 tests~: quelle p-value doit avoir la vingtième plus petite pour que vingt hypothèses soient rejetées~?
\end{exercise}

\begin{exercise}[$\star$]\label{exo:s1:pairs-and-baskets:3}
Une paire est ouverte puis clôturée une fois, à cinq points de base par unité échangée, avec 1~\$ par jambe. Combien coûte l'aller-retour, en points de base d'une jambe~?
\end{exercise}

\begin{exercise}[$\star\star$]\label{exo:s1:pairs-and-baskets:4}
L'écart-type de formation de l'exemple d'ouverture est de 2.4~\%. À quel écart la paire s'est-elle ouverte, et pourquoi l'offre publique d'achat lui a-t-elle fait perdre bien davantage que cela~?
\end{exercise}

\begin{exercise}[$\star\star$]\label{exo:s1:pairs-and-baskets:5}
Pourquoi la méthode de la distance favorise-t-elle les paires d'actions à faible volatilité, et pourquoi cela compte-t-il lorsque le marché n'a pas de substituts~?
\end{exercise}

\begin{exercise}[$\star\star$]\label{exo:s1:pairs-and-baskets:6}
Expliquer comment 97~\% des découvertes de Benjamini--Hochberg peuvent être fausses lorsque la procédure est réglée à 5~\%.
\end{exercise}

\begin{exercise}[$\star\star\star$]\label{exo:s1:pairs-and-baskets:7}
\emph{Programmation.} À l'aide de \texttt{discoveries(0.1)}, calculer la part des paires implantées parmi les paires à nouveau rejetées l'année suivante, et la comparer
à leur part parmi les rejets de l'année de formation.
\end{exercise}

\begin{exercise}[$\star\star\star$]\label{exo:s1:pairs-and-baskets:8}
\emph{Trouver la faille.} «~Nous avons testé les 50\,000 paires intra-sectorielles, retenu celles significatives à 1~\% après correction des fausses découvertes, puis backtesté les
survivantes~: ratio de Sharpe de 3.~»
\end{exercise}

\section{Problème~: La paire qui ne s'est jamais refermée}

\begin{problem}[{Problème du week-end --- la sélection est la stratégie}]\label{pb:s1:pairs-and-baskets:1}
Les cinq méthodes de sélection du chapitre sur \texttt{firm.synthmkt}, avec et sans substituts implantés.

\medskip\noindent\textbf{Partie I --- Les méthodes.}
\begin{enumerate}
\item Définir le \omterm{def:s1:pairs-and-baskets:pairs}{pairs trading} ainsi que les \omterm{def:s1:pairs-and-baskets:distance}{périodes de formation} et de trading de la méthode de la distance.
\item Énoncer la règle de Gatev, Goetzmann et Rouwenhorst et leur résultat principal.
\item Comment les méthodes de cointégration et de copule sélectionnent-elles les paires et émettent-elles leurs signaux~?
\item Qu'ont trouvé Rad, Low et Faff, avant et après les coûts~?
\end{enumerate}

\medskip\noindent\textbf{Partie II --- Le test synthétique.}
\begin{enumerate}[resume]
\item Qu'implante \texttt{twin\_share}, et pourquoi le chapitre en a-t-il besoin~?
\item Combien de paires implantées chaque méthode sélectionne-t-elle~?
\item Donner les ratios de Sharpe nets sur les deux marchés.
\item Décrire l'exemple d'ouverture et expliquer pourquoi la paire ne s'est jamais refermée.
\end{enumerate}

\medskip\noindent\textbf{Partie III --- Tester à grande échelle.}
\begin{enumerate}[resume]
\item Combien de tests le chapitre effectue-t-il, et quel est leur degré de calibration~?
\item Quelle part des découvertes de Benjamini--Hochberg est fausse, et pourquoi~?
\item Que fait la méthode en deux étapes, et pourquoi ne bat-elle pas la distance~?
\item Pourquoi le \omterm{def:s1:pairs-and-baskets:twin}{jumeau synthétique} perd-il~?
\end{enumerate}

\medskip\noindent\textbf{Partie IV --- Le verdict.}
\begin{enumerate}[resume]
\item Énoncer le \emph{résultat nommé}~: la part des paires cointégrées pendant l'année de formation qui restent cointégrées pendant l'année de trading, et le
ratio de Sharpe net du portefeuille de paires après coûts.
\item Que révèle la persistance sur les vraies et les fausses paires~?
\item Qu'ont trouvé Do et Faff sur la tendance et sur les crises~?
\item Qu'est-ce qui a remplacé la paire classique~?
\item Comment vérifier la calibration d'un test avant de l'utiliser à grande échelle~?
\item Que faut-il exiger d'une paire avant de la trader~?
\item Comment protéger un portefeuille de paires contre les offres publiques d'achat~?
\item En une phrase~: que demande une stratégie de paires à son marché~?
\end{enumerate}
\end{problem}

\section{Questions d'entretien}

\begin{interviewq}[$\star$ \iqroles{chercheur, trader}]\label{iq:s1:pairs-and-baskets:1}
Comment choisiriez-vous des paires pour une stratégie de \omterm{def:s1:pairs-and-baskets:pairs}{pairs trading}~?
\end{interviewq}

\begin{interviewq}[$\star\star$ \iqroles{chercheur}]\label{iq:s1:pairs-and-baskets:2}
Corrélation ou cointégration pour les paires~: quelle est la différence, et laquelle compte~?
\end{interviewq}

\begin{interviewq}[$\star\star$ \iqroles{chercheur}]\label{iq:s1:pairs-and-baskets:3}
Vous testez 100\,000 paires pour la cointégration. Comment maîtrisez-vous les fausses découvertes, et que peut-il encore mal tourner~?
\end{interviewq}

\begin{interviewq}[$\star\star$ \iqroles{risque, trader}]\label{iq:s1:pairs-and-baskets:4}
L'une des jambes d'une paire ouverte fait l'objet d'une offre publique d'achat. Que faites-vous~?
\end{interviewq}

\begin{interviewq}[$\star\star$ \iqroles{développeur}]\label{iq:s1:pairs-and-baskets:5}
Calculez efficacement la somme des carrés des distances entre les prix normalisés de toutes les paires de 3\,000 actions.
\end{interviewq}

\begin{interviewq}[$\star\star\star$ \iqroles{chercheur}]\label{iq:s1:pairs-and-baskets:6}
Un écart suit un AR(1) de coefficient $\phi$ et d'écart-type stationnaire $\sigma$. Quel est le temps moyen pour revenir à mi-chemin, et comment
la puissance du test de Dickey--Fuller dépend-elle de $\sigma$~?
\end{interviewq}
